\section{Root independence, uniqueness, and the comparison theorem}\label{sec:comparison}

\subsection{Root independence of the tree coefficient}

\begin{lemma}[triangle and bow-tie values]\label{lem:A-small-values}
\begin{enumerate}
\item[(i)] For every labelled triangle satisfying (G1), all turns
have a common sign $\tau$, and $A_g=-\tau$ at every physical root.
\item[(ii)] For the bow-tie of Definition~\ref{def:star},
$A_g(\sigma^kK_0)=-1$ for every physical root $g$ and every integer $k$.
\end{enumerate}
\status{new (round~4; extracted from the SM4 proof of Theorem~\ref{thm:root-indep-proof}; adopted from the Codex proposal sm4\_round4\_exports\_v1, P-25-24, F117-SMALL-LEMMA, re-read by the author; F-25-117; cleared by bench~A on SM5 (2026-09-06T00:32:58Z))}
\end{lemma}
\begin{proof}
For a noncollinear triangle the three turn triples are cyclic
permutations of one ordered triple. Lemma~\ref{lem:chi-basic}(i)
therefore gives the common sign $\tau$. For any physical root,
the boundary word has two leaves. Its unary root contribution is
$b_{[0,2]}=0$, since its only ordinary composition has two parts
and vanishes by Lemma~\ref{lem:gates-nonzero}. The binary root
has $d=h=\chi(a_2,a_1,a_0)=-\tau$, so its factor
$(d+h)/2$ is $-\tau$. This proves (i) directly on (G1).

For (ii), use the four coordinates of Definition~\ref{def:star}.
All two-leaf open sums vanish, because their near
and far signs agree. On a boundary word $(a_0,a_1,a_2,a_3)$, put
$d_1=\chi(a_2,a_1,a_0)$, $d_2=\chi(a_3,a_2,a_1)$,
$h_1=\chi(a_3,a_1,a_0)$ and $h_2=\chi(a_3,a_2,a_0)$.
Every two-part composition has a two-leaf child and vanishes.
The unary root term is therefore minus the ordinary three-part gate
product, and the three-part root term is its barred counterpart. Hence
\begin{equation}\label{eq:root-bowtie-recursion}
A_g=\frac{(d_1+h_1)(d_2+h_2)-(d_1-h_1)(d_2-h_2)}4.
\end{equation}
Expanding the two products cancels $d_1d_2$ and $h_1h_2$ and leaves
$2d_1h_2+2h_1d_2$. Thus $A_g=(d_1h_2+h_1d_2)/2$.
At the four physical roots of $K_0$, direct determinants of its displayed
coordinates give
\begin{equation}\label{eq:root-bowtie-signs}
\begin{array}{c|rrrr|r}
g&d_1&d_2&h_1&h_2&A_g\\\hline
1&-1&1&-1&1&-1\\
2&1&1&-1&-1&-1\\
3&1&-1&1&-1&-1\\
4&-1&-1&1&1&-1
\end{array}
\end{equation}
Cyclic shifts merely permute these four root words by
Proposition~\ref{prop:A-reversal}(i). Consequently every root of every
$\sigma^kK_0$ has value $-1$.
\end{proof}

\begin{theorem}[root independence]\label{thm:root-indep-proof}
For every generic polygon $P$ and any two of its edges $g,h$, $A_g(P)=A_h(P)$.
\status{new (round~4: the anchor parameter chosen below the geometric bound and the two root-specific soft bounds, F-25-119; the triangle value and the bow-tie table extracted to Lemma~\ref{lem:A-small-values} and the negative stars handled by reversal, F-25-117; Codex proposals P-25-23 and P-25-24 adopted or adapted)}
\end{theorem}
\begin{proof}
Strong induction on $n$. The case $n=3$ follows from
Lemma~\ref{lem:A-small-values}(i).

Let $n\geq4$, let $P$ be generic of rotation $r$, and let $g\neq h$ be two
edge labels. For a generic labelled tuple $Q$ with $n$ vertices put
$\Delta_{gh}(Q)=A_g(Q)-A_h(Q)$. Choose a target $Z$ in the fibre $(n,r)$ as
follows.
\begin{itemize}
\item If $(n,r)$ is not minimal: $Z$ is a zero anchor whose insertion vertex
$j$ is chosen with $j\notin\{g,h\}$ (Proposition~\ref{prop:anchors-exist};
there are $n-1\geq3$ choices of $j$). Its soft edge is $E_j(Z)$.
Choose its parameter below the geometric bound of
Definition~\ref{def:anchors} and below the two bounds in
Theorem~\ref{thm:A-soft} for the roots $g$ and $h$.
All three bounds are positive, so such a choice exists.
\item If $(n,r)$ is minimal with $|r|\geq2$: $Z=K_r$.
\item If $(n,r)=(4,0)$: $Z=K_0$.
\end{itemize}
Lemma~\ref{lem:transport} gives $k$ ($k=0$ in the first two cases) and a path
from $P$ to $\sigma^kZ$ with finitely many simple walls. Along the path
$\Delta_{gh}$ is constant: on generic segments by
Proposition~\ref{prop:A-chamber}; at (T), (E), (C) by
Theorem~\ref{thm:A-R3E}, which gives zero jump for each root separately; at a
flat wall at $i$ the jump of $\Delta_{gh}$ is
$A_{D_i(g)}(Q)-A_{D_i(h)}(Q)$ with $Q$ the generic deletion
(Theorem~\ref{thm:A-S3}), which vanishes by induction; at a vertex--edge wall
the jump is
$s\bigl(A_{h_1}(\lambda_1)A_{h_2}(\lambda_2)-A_{h_1'}(\lambda_1)A_{h_2'}(\lambda_2)\bigr)$
with $(h_1,h_2)=H(g)$, $(h_1',h_2')=H(h)$ and generic halves
(Theorem~\ref{thm:A-S7}), which vanishes by induction applied to each half.
Hence $\Delta_{gh}(P)=\Delta_{gh}(\sigma^kZ)$.

At the target: in the non-minimal case $k=0$ and neither $E_g$ nor $E_h$ is
the soft edge of $Z$, so Theorem~\ref{thm:A-soft} in the mixed sector gives
$A_g(Z)=A_h(Z)=0$. For $Z=K_r$ with $r\geq2$: by Lemma~\ref{lem:star-generic}(i) and
Proposition~\ref{prop:A-reversal}(i),
$A_{g+1}(K_r)=A_g(\sigma K_r)=A_g(R(K_r))=A_g(K_r)$, the last step because
every chirotope, hence every gate, is invariant under a rotation of the plane;
so all roots of $K_r$ have one value. For $r\leq-2$, $K_r=\overline{K_{|r|}}$
by Definition~\ref{def:star}, and Proposition~\ref{prop:A-reversal}(ii) gives
$A_h(K_r)=(-1)^NA_{1-h}(K_{|r|})=-A_{1-h}(K_{|r|})$ with $N=2|r|+1$ odd; since
$h\mapsto1-h$ is a bijection of $\ZZ/N$, all roots of $K_r$ have one value as
well. Cyclic covariance, Proposition~\ref{prop:A-reversal}(i), covers every
shift. For $Z=K_0$, Lemma~\ref{lem:A-small-values}(ii) gives $-1$ at every
root of every shift. In every target case $\Delta_{gh}(\sigma^kZ)=0$.
\end{proof}

\begin{corollary}[$A$ is a chamber function on polygons]\label{cor:A-lawful}
$A(P):=A_g(P)$ (any $g$) is well defined on generic polygons and satisfies:
chamber constancy and silence (Proposition~\ref{prop:A-chamber},
Theorem~\ref{thm:A-R3E}(ii)); the flat law
$A(P_{\rm right})-A(P_{\rm left})=A(P(0)\setminus j)$; the cusp law
$A(P_{\rm loop})-A(P_{\rm no})=-\kappa A(P(0)\setminus j)$ when the deletion
satisfies (G1); the vertex--edge law $A(P_+)-A(P_-)=sA(\lambda_1)A(\lambda_2)$;
the triple law $A(P_+)=A(P_-)$; the soft theorem
$A(P_\varepsilon)=\frac{\chi_-+\chi_+}2A(P)$ in every sector; reversal
$A(\overline P)=(-1)^nA(P)$; and $A(K_1)=-1$, $A(K_{-1})=+1$ on the
counterclockwise and clockwise triangles. \status{new (round~4 assembly: Theorem~\ref{thm:root-indep-proof} (new) removes the root choice; Proposition~\ref{prop:A-reversal}(i),(ii), Theorems~\ref{thm:A-S3}, \ref{thm:A-S4}, \ref{thm:A-S7}, \ref{thm:A-R3E} and~\ref{thm:A-soft} refereed; Proposition~\ref{prop:A-chamber} proved; Lemma~\ref{lem:A-small-values} (new) supplies the triangle values; F-25-117)}
\end{corollary}
\begin{proof}
Theorem~\ref{thm:root-indep-proof} removes the root choice, and
Proposition~\ref{prop:A-reversal}(i) gives cyclic descent.
Theorems~\ref{thm:A-S3}, \ref{thm:A-S4}, \ref{thm:A-S7}, \ref{thm:A-R3E},
\ref{thm:A-soft} (evaluated at any root other than the soft edge),
Proposition~\ref{prop:A-reversal}(ii), and the triangle values in
Lemma~\ref{lem:A-small-values}(i) give the listed laws and base values.
\end{proof}

\begin{corollary}[root independence and cyclic descent of the continuation]\label{cor:weak-root-indep}
For $n\geq3$ and any two roots $g,h$, the continuations of
Theorem~\ref{thm:A-continuation} satisfy $\widetilde A_g=\widetilde A_h$ on
$\mathcal W_n$. For every root $g$ and every labelled weakly generic tuple
$P$,
\begin{equation}\label{eq:weak-root-covariance}
\widetilde A_g(\sigma P)=\widetilde A_{g+1}(P),
\end{equation}
with root indices read modulo $n$; consequently
$\widetilde A_g(\sigma P)=\widetilde A_g(P)$, and each $\widetilde A_g$ is
a function on weakly generic polygons.
\status{new (round~3; adapted from the Codex proposal cor:weak-root-indep, P-25-17, lane sm2\_integrated\_review\_v1, with bench~A's remark of 2026-09-05T18:31:35Z; no Hypothesis~R; F-25-68, F-25-41)}
\end{corollary}
\begin{proof}
By Theorem~\ref{thm:root-indep-proof}, $A_g=A_h$ on $\mathcal U_n$. By
Theorem~\ref{thm:A-continuation}, $\widetilde A_h$ is constant on visible
chambers and agrees with $A_h=A_g$ on $\mathcal U_n$; the uniqueness clause
of that theorem for the root $g$ gives $\widetilde A_h=\widetilde A_g$.

The cyclic shift $\sigma$ permutes the coordinates of a labelled tuple, and
every condition of Definition~\ref{def:weak} is a condition on the polygon:
nonzero edge vectors and turns, no vertex on a nonincident closed edge,
remote edges meeting at most once and transversally, and (G2) are all
invariant under relabelling. Hence $\sigma$ and $\sigma^{-1}$ map
$\mathcal W_n$ onto itself, are continuous, and permute its connected
components, the visible chambers; likewise for $\mathcal U_n$. The
function $P\mapsto\widetilde A_g(\sigma P)$ is therefore constant on every
visible chamber, and at a generic tuple its value is
$A_g(\sigma P)=A_{g+1}(P)$ by Proposition~\ref{prop:A-reversal}(i). The
uniqueness clause of Theorem~\ref{thm:A-continuation} for the root $g+1$
gives $\widetilde A_g\circ\sigma=\widetilde A_{g+1}$, which
is~\eqref{eq:weak-root-covariance}. With the first assertion,
$\widetilde A_g(\sigma P)=\widetilde A_{g+1}(P)=\widetilde A_g(P)$;
iterating, $\widetilde A_g$ is invariant under every power of $\sigma$ and
descends to the orbits of Definition~\ref{def:polygon}.
\end{proof}

\begin{definition}[the amplitude on the weakly generic locus]\label{def:weak-amplitude}
For a weakly generic polygon $P\in\mathcal W_n$ put
\begin{equation}\label{eq:weak-amplitude}
A^{\rm cont}(P)=\widetilde A_g(P)\qquad(\text{any root }g),
\end{equation}
which is well defined on polygons by Corollary~\ref{cor:weak-root-indep}.
On $\mathcal U_n$ it is the function $A(P)=A_g(P)$ of
Corollary~\ref{cor:A-lawful}, by the agreement clause of
Theorem~\ref{thm:A-continuation}. The main text's amplitude on its generic
locus, which is $\mathcal W_n$ (Definition~\ref{def:weak}), is read as
$A^{\rm cont}$. On the silent locus $\mathcal S_n$ this is a definition,
not an identity with the tree sum of Definition~\ref{def:treesum}, which
assigns no value to a gate at zero (Remark~\ref{rem:silent-amplitude}).
\end{definition}

\begin{corollary}[finite evaluation of the continuation]\label{cor:weak-amplitude}
For every $P\in\mathcal W_n$ and every root $g$,
\begin{equation}\label{eq:weak-amplitude-polynomial}
A^{\rm cont}(P)=\mathsf A_g(\chi(P)),
\end{equation}
the polynomial form of Corollary~\ref{cor:polyform} evaluated at the
chirotope of $P$, zero entries included. Moreover $A^{\rm cont}$ is the
unique function on $\mathcal W_n$ that is constant on visible chambers and
agrees with $A$ on $\mathcal U_n$.
\status{new (round~3; adapted from the Codex proposal cor:weak-amplitude, F68-P2, P-25-17, with bench~A's correction of 2026-09-05T18:31:35Z; uses Corollary~\ref{cor:polyform}, tag \textsc{new}, whose two lemmas clear bench~A's referee, R-25-21; no Hypothesis~R; F-25-68); assembly: Definition~\ref{def:weak-amplitude}; Corollary~\ref{cor:polyform} new (its two lemmas refereed, R-25-21; its round-4 lemmas unrefereed); Theorem~\ref{thm:A-continuation} new; Corollary~\ref{cor:A-lawful} new (round~4)}
\end{corollary}
\begin{proof}
By Definition~\ref{def:weak-amplitude}, $A^{\rm cont}(P)=\widetilde A_g(P)$,
and Corollary~\ref{cor:polyform} identifies this with $\mathsf A_g(\chi(P))$
on all of $\mathcal W_n$. For the uniqueness, $A^{\rm cont}$ is constant on
visible chambers and agrees with $A=A_g$ on $\mathcal U_n$
(Theorem~\ref{thm:A-continuation}, Corollary~\ref{cor:A-lawful}); any
function with these two properties equals $\widetilde A_g=A^{\rm cont}$ by
the uniqueness clause of Theorem~\ref{thm:A-continuation}.
\end{proof}

\subsection{Uniqueness}

\begin{theorem}[uniqueness; Theorem~1 of the main text]\label{thm:uniqueness}
Let $F$ be a function assigning an integer to every generic polygon of every
arity $n\geq3$ and satisfying:
\begin{enumerate}
\item[(a)] $F$ is constant on chambers and unchanged across simple silent
walls (types (E) and (C));
\item[(b)] the flat law $F(P_{\rm right})-F(P_{\rm left})=F(P(0)\setminus j)$
at every simple flat wall;
\item[(c)] the vertex--edge law $F(P_+)-F(P_-)=s\,F(\lambda_1)F(\lambda_2)$
at every simple vertex--edge wall, both branches;
\item[(d)] $F(P_+)=F(P_-)$ at every simple triple wall;
\item[(e)] the soft theorem $F(P_\varepsilon)=\frac{\chi_-+\chi_+}2F(P)$ for
all small $\varepsilon$, at every soft insertion of an admissible vector into
a generic polygon;
\item[(f)] $F(K_1)=-1$ and $F(K_{-1})=+1$.
\end{enumerate}
Then $F(P)=A(P)$ for every generic polygon $P$. \status{new (round~4: the anchor step cites hypotheses (a) and (e), as Proposition~\ref{prop:anchor-values} now requires, F-25-119)}
\end{theorem}
\begin{proof}
$A$ satisfies (a)--(f) by Corollary~\ref{cor:A-lawful}. Put $\Delta=F-A$, a
function on generic polygons, and argue by strong induction on $n$.

\emph{Base $n=3$.} The generic triangles form two chambers, the
counterclockwise and the clockwise triangles, containing $K_1$ and $K_{-1}$
respectively; by (a) and (f), $\Delta=0$ on both.

\emph{Step.} Let $n\geq4$ and assume $\Delta=0$ on every generic polygon with
fewer than $n$ vertices. Let $P$ be generic with $\rot(P)=r$, so that $(n,r)$
is admissible, and let $Z$ be an anchor for $(n,r)$
(Definition~\ref{def:anchors}, Proposition~\ref{prop:anchors-exist}).
Lemma~\ref{lem:transport} gives a path from $P$ to $\sigma^kZ$, which as a
path of polygons ends at $Z$, with finitely many simple walls of types (F),
(V), (T), (E), (C). Between walls $\Delta$ is constant by (a); at (E), (C) it
is unchanged by (a); at (T) by (d). At a flat wall at $j$ the jump of
$\Delta$ is $\Delta(P(0)\setminus j)=0$ by induction, the deletion being
generic with $n-1$ vertices. At a vertex--edge wall the jump is
$s\bigl(F(\lambda_1)F(\lambda_2)-A(\lambda_1)A(\lambda_2)\bigr)
=s\bigl(\Delta(\lambda_1)F(\lambda_2)+A(\lambda_1)\Delta(\lambda_2)\bigr)=0$
by induction, the halves being generic with fewer than $n$ vertices. Hence
$\Delta(P)=\Delta(Z)$. At the anchor, Proposition~\ref{prop:anchor-values}
applies to $F$ and to $A$ (hypotheses (a) and (e),
Corollary~\ref{cor:A-lawful}):
for a zero anchor both values are $0$; for a loop anchor with parent $P'$
(a generic polygon with $n-1$ vertices) and insertion vertex $j$,
$F(Z)=\tau_j(P')F(P')$ and $A(Z)=\tau_j(P')A(P')$, and $\Delta(P')=0$ by
induction. So $\Delta(Z)=0$ and $\Delta(P)=0$.
\end{proof}

\begin{remark}[what the theorem does not assume]
No invariance under reversal is assumed. Cyclic relabelling is already
built into $F$ being a function on polygons, and the two triangle values
are prescribed separately. The cusp law is not assumed either; transport
paths avoid cusps (Theorem~\ref{thm:relgp}). Anchor values follow from the
soft theorem together with chamber constancy on the initial soft family.
The prescribed values at the triangle stars $K_1$ and $K_{-1}$ are the
base data; no value at a higher star or at the bow-tie is assumed.
Those remaining values in Section~\ref{sec:anchors} are consequences.
\end{remark}

\begin{corollary}[uniqueness on the weakly generic locus]\label{cor:uniqueness-weak}
For every $n\geq3$ let $F$ assign an integer to every weakly generic polygon
in $\mathcal W_n$ (Definition~\ref{def:weak}). Suppose that $F$ is constant on
every visible chamber and that its restrictions to the generic loci
$\mathcal U_n$, for all arities, satisfy (a)--(f) of
Theorem~\ref{thm:uniqueness}. Then $F(P)=\widetilde A_g(P)$ for every
$P\in\mathcal W_n$ and every root $g$; in particular the continuation
$\widetilde A_g$ of Theorem~\ref{thm:A-continuation} does not depend on $g$.
\status{new (author, round 1; adapted from bench A P-25-4 item 13 and Codex P-25-9; F-25-41); assembly: Theorems~\ref{thm:uniqueness}, \ref{thm:root-indep-proof} and~\ref{thm:A-continuation} new (bench~A cleared the first two on their own arguments on SM3, the third read in full on SM2)}
\end{corollary}
\begin{proof}
By Theorem~\ref{thm:uniqueness}, $F=A$ on every $\mathcal U_n$, and by
Theorem~\ref{thm:root-indep-proof}, $A=A_g$ there for every root $g$. Fix a
root $g$. By Theorem~\ref{thm:A-continuation}, $\widetilde A_g$ is the unique
function on $\mathcal W_n$ that is constant on visible chambers and agrees
with $A_g$ on $\mathcal U_n$. The function $F$ is constant on visible chambers
by hypothesis and agrees with $A_g$ on $\mathcal U_n$ by the first sentence;
hence $F=\widetilde A_g$ on $\mathcal W_n$. The left side does not involve
$g$, so $\widetilde A_g=\widetilde A_h$ for any two roots $g,h$.
\end{proof}

\begin{remark}[the weak-locus hypothesis; Theorem~1 of the main text]
The main text's Theorem~1 concludes on its generic polygons, which are the
weakly generic polygons of Definition~\ref{def:weak}, and its
chamber-constancy hypothesis is constancy on the visible chambers. Taking
any root $g$ in Corollary~\ref{cor:uniqueness-weak} and
Definition~\ref{def:weak-amplitude} gives $F=A^{\rm cont}$ on
$\mathcal W_n$: the corollary is the form in which
Theorem~\ref{thm:uniqueness} covers the main text's statement, its
amplitude at a silent configuration being read as the continuation, which
is how this document defines it (Remark~\ref{rem:silent-amplitude}); no
gate is evaluated at zero. Constancy on the chambers of $\mathcal U_n$
alone would leave the values of $F$ on the silent locus unspecified. The
common value is also the polynomial value $\mathsf A_g(\chi(P))$ by
Corollary~\ref{cor:weak-amplitude}, which is not used in the proof.
\end{remark}

\subsection{The comparison theorem}

\begin{theorem}[comparison]\label{thm:comparison}
Assume Hypothesis~R. Then $C(P)=A(P)$ for every generic polygon $P$.
\status{new (assembly; Proposition~\ref{prop:C-silent}, Theorems~\ref{thm:C-S3} and~\ref{thm:C-S7} refereed; Proposition~\ref{prop:C-chamber} and Lemma~\ref{lem:corner-values} held; Theorem~\ref{thm:C-soft} new (round~4), unrefereed; Theorem~\ref{thm:uniqueness} new; conditional on Hypothesis~R; round~5: hypothesis (d) discharged by citing Hypothesis~\ref{hyp:R} by label, so that the ordering instrument records the consumption, F-25-160)}
\end{theorem}
\begin{proof}
$C$ satisfies (a)--(f) of Theorem~\ref{thm:uniqueness}: (a)
Propositions~\ref{prop:C-chamber}, \ref{prop:C-silent}; (b)
Theorem~\ref{thm:C-S3}; (c) Theorem~\ref{thm:C-S7}; (d) Hypothesis~\ref{hyp:R}; (e)
Theorem~\ref{thm:C-soft}; (f) Lemma~\ref{lem:corner-values}(i)
gives corner coefficient $1$ for either oriented triangle. Its only
decomposition is empty, so Definition~\ref{def:C} gives $-1$ when
$\ell=3$ and $+1$ when $\ell=0$, as required.
\end{proof}

\begin{corollary}[laws inherited by $C$]\label{cor:C-inherits}
Under Hypothesis~R, $C$ satisfies every identity of
Corollary~\ref{cor:A-lawful} on the domains stated there, in particular
the cusp law
$C(P_{\rm loop})-C(P_{\rm no})=-\kappa\,C(P(0)\setminus j)$.
\status{new (printed in round~1, C013; restricted to Corollary~\ref{cor:A-lawful}, F-25-11; round~3: proof completed with the domain checks on bench~A's named step, adapted from the Codex proposal P-25-19, F-25-11)}
\end{corollary}
\begin{proof}
Theorem~\ref{thm:comparison} identifies $C$ with $A$ on every generic
polygon, and Definition~\ref{def:C} defines $C$ on generic polygons only;
so before $C=A$ is substituted in an identity of
Corollary~\ref{cor:A-lawful}, every argument of that identity must be
shown generic. The two sides of a wall germ are generic by
Definition~\ref{def:germ}. At a simple flat wall the deletion
$P(0)\setminus j$ is generic by Lemma~\ref{lem:children}(i), and at a
simple vertex--edge wall both halves are generic by
Lemma~\ref{lem:children}(ii). The soft family $P_\varepsilon$ is generic
for all sufficiently small $\varepsilon>0$ in every sector by
Lemma~\ref{lem:soft-generic}. Reversal permutes the vertex triples and
preserves the edge segments, so it preserves (G1) and (G2). The two
triangles $K_{\pm1}$ are generic by Lemma~\ref{lem:star-generic}(iii).

For the cusp law let the simple cusp wall be at $j$ and write
$A=\mu_{j-1}(0)$, $M=\mu_j(0)$ and $B=\mu_{j+1}(0)$. The deletion
$Q=P(0)\setminus j$ has the vertices $\mu_i(0)$, $i\neq j$, and its edges
are the parent edges $E_i(0)$, $i\notin\{j-1,j\}$, together with the fused
edge $[A,B]$; it satisfies (G1) by the hypothesis of the cusp law. For (G2):
by Definition~\ref{def:walls}(K) the three points are collinear with $M$
outside the closed segment $[A,B]$, so $[A,B]$ is contained in the longer of
$[A,M]=E_{j-1}(0)$ and $[M,B]=E_j(0)$, and its relative interior in the
relative interior of that edge. Suppose three distinct edges of $Q$ had a
common point in their relative interiors. Two adjacent edges of $Q$ are not
collinear, by (G1) for $Q$, and so meet only at their common vertex; hence
the three edges are pairwise remote in $Q$. Send the fused edge, if it is
among them, to the longer parent edge containing it, and every other edge
to itself. This map is injective, since neither $E_{j-1}(0)$ nor $E_j(0)$
is an edge of $Q$, and the common point lies in the relative interiors of
the three image edges. The image edges are pairwise remote in $P(0)$: two
surviving edges are adjacent in $P(0)$ exactly when they share a vertex,
which is a vertex of $Q$, hence exactly when they are adjacent in $Q$; and
the parent neighbours of the longer cusp edge are the other cusp edge,
which is not an edge of $Q$, and one of the two edges of $Q$ adjacent to
the fused edge, both excluded because the three edges are pairwise remote
in $Q$. This is a point in the relative interiors of three pairwise remote
edges of the centre, contrary to $Z_{\rm c}=\varnothing$ in
Definition~\ref{def:walls}(K). Thus $Q$ is generic; the argument does not
require the cusp to be empty.

Now apply each identity of Corollary~\ref{cor:A-lawful} to its stated
arguments and substitute $C=A$ at every one of them, including the
deletions and halves just checked. For a chamber-side identity each
sufficiently small punctured side lies in one generic chamber, on which
$C=A$ holds pointwise, so the constant side values of $C$ and of $A$
agree; no value at the wall centre and no limit is used. This gives every
identity for $C$ on the domain stated in Corollary~\ref{cor:A-lawful},
including the displayed cusp law.
\end{proof}

\begin{corollary}[comparison on the weakly generic locus]\label{cor:comparison-weak}
Assume Hypothesis~R. Then $C(P)=A^{\rm cont}(P)$ for every weakly generic
polygon $P\in\mathcal W_n$, with $C$ the state sum of
Definition~\ref{def:C-weak} and $A^{\rm cont}$ the continuation of
Definition~\ref{def:weak-amplitude}.
\status{new (round~3; adapted from the Codex proposal W-P3, P-25-13, lane sm2\_weak\_c\_v1; conditional on Hypothesis~R; F-25-90); assembly: Theorem~\ref{thm:comparison} new, conditional; Theorem~\ref{thm:root-indep-proof} new; Lemma~\ref{lem:C-weak-local} new (cleared by bench~A on SM4, R-25-36); Theorem~\ref{thm:A-continuation} new}
\end{corollary}
\begin{proof}
On $\mathcal U_n$, Theorem~\ref{thm:comparison} gives $C=A$, and $A=A_g$
for every root $g$ by Theorem~\ref{thm:root-indep-proof}. By
Lemma~\ref{lem:C-weak-local}, $C$ is constant on every visible chamber of
$\mathcal W_n$. Fix a root $g$. The uniqueness clause of
Theorem~\ref{thm:A-continuation} identifies every function on
$\mathcal W_n$ that is constant on visible chambers and agrees with $A_g$
on $\mathcal U_n$ with $\widetilde A_g=A^{\rm cont}$; hence
$C=A^{\rm cont}$. No identification of a value at a silent configuration
with the tree sum of Definition~\ref{def:treesum} is used.
\end{proof}

\begin{remark}[Dependence on Hypothesis R]\label{rem:conditional}
The printed rooted tree laws, root-independence, uniqueness and the
continuation (Corollaries~\ref{cor:weak-root-indep}
and~\ref{cor:weak-amplitude}) do not use Hypothesis~\ref{hyp:R}.
Theorem~\ref{thm:comparison}, Corollary~\ref{cor:C-inherits} and
Corollary~\ref{cor:comparison-weak} do use it. A later transfer from a
value of $C$ to a value of $A$ through that comparison has the same
condition. The independent source premises in the registry remain
explicit. The polynomial identification of Corollary~\ref{cor:polyform} on
all of $\mathcal W_n$ and the linear relations of
Section~\ref{sec:relations} are printed with the tag \textsc{new}; their
proofs were assembled in round~1 of Phase~2.5, the two lemmas of the former
were refereed by bench~A on SM2 (ledger row F-25-25, closed by ruling
R-25-21), and the exception request of decision item D3 for the latter was
withdrawn (ruling R-25-17: the former step~(ii) is obviated, not proved).
\end{remark}
